\chapter{Pay Regulation and Tax by Location}\label{ch:in:pay-regulation-and-tax-by-location}

The European Union's 2013 capital directive capped a bank risk-taker's variable pay at the level of fixed pay, or twice
it with the shareholders' approval; the United Kingdom removed the cap from 31 October 2023, and the cap has not applied
to EU investment firms since 2021. Where the cap binds, pay is shifted from bonus to base; where it does not, the bonus
can be several times the base, as chapter~14's high-earner figures showed. And once a package is set, where the employee
lives decides how much of it is left: a package of one million dollars leaves about \$516\,000 in Amsterdam and all of it
in Dubai, and a special regime for incomers can move the answer by tens of thousands. This chapter covers the rules that
shape pay and the taxes that take from it, for eight places where the industry employs people, and computes what a
package is worth after tax in each. It explains how the systems work; it is not tax advice, and a reader's own case
turns on facts the chapter does not see.

\section{Pay regulation: the bonus cap, deferral rules and who they cover}

\begin{definition}[Bonus cap]\label{def:in:pay-regulation-and-tax-by-location:cap}
A \emph{bonus cap}\index{bonus cap} is a legal limit on the ratio of variable to fixed pay. In the EU's capital
requirements directive the variable component ``shall not exceed 100\,\% of the fixed component'' for each individual
whose work has a material impact on the institution's risk, with up to 200\,\% allowed if the owners approve.
\end{definition}

Three facts decide whether a rule applies to a given person.
\begin{itemize}
\item \textbf{The employer's status.} The cap and the deferral rules of chapter~13 apply to credit institutions and, in
the UK, to banks, building societies and the investment firms the regulator designates; the EU's investment firms have
been under their own regime since 2021, without the cap.
\item \textbf{The person's role.} Within a covered firm the strictest rules apply to material risk takers (Book~16,
chapter~10), identified by their role and their pay.
\item \textbf{The place.} The UK's removal of the cap in 2023 and its 2025 changes to deferral apply to UK-regulated firms;
which regulator's rules apply to a person depends on the entity that employs them and where it is authorised.
\end{itemize}
The cap does not limit total pay: a firm that cannot pay a larger bonus can pay a larger base, which is what the UK
regulators gave as the cap's effect when they removed it (Book~16, chapter~10). For a candidate, the cap changes the mix
more than the level: more is fixed, which is safer, and less can be clawed back or deferred. US rules take a different
route: listed issuers must claw back incentive pay after an accounting restatement (Book~16, chapter~10).

\begin{example}[The cap and the mix]\label{ex:in:pay-regulation-and-tax-by-location:mix}
A bank wants to offer a risk-taker a total of \euro1\,000\,000 in a normal year. Under the EU cap without the owners'
approval, variable pay may not exceed fixed pay, so fixed pay must be at least \euro500\,000; with approval (200\%), at
least \euro333\,333; uncapped, as in the UK since 2023, the bank may offer, say, \euro300\,000 fixed and \euro700\,000
variable. In a year when the variable award is zero, the employee is paid \euro500\,000, \euro333\,333 or \euro300\,000:
the cap is insurance for the employee and a fixed cost for the bank, which is why the bank's own view of it, in Book~16,
chapter~10, is about the loss-absorbing share of pay.
\end{example}

\section{Income tax on salary and bonus}

\begin{definition}[Marginal tax rate, average tax rate]\label{def:in:pay-regulation-and-tax-by-location:rates}
The \emph{marginal tax rate}\index{marginal tax rate} is the share of the next unit of pay taken in income tax and
employee social contributions; the \emph{average tax rate}\index{average tax rate} is the share of the whole of pay so
taken. A bonus is taxed at the marginal rates it falls into; a comparison of packages uses the average rate.
\end{definition}

In every system the chapter models, a bonus is taxed as employment income like salary: withholding on a bonus may follow
special tables during the year, but the year's tax is the same whether pay comes as salary or bonus. The systems differ in their
rates, their thresholds, their social contributions and whether contributions are capped.

\begin{dated}{2026-09}{Rates for a single employee, 2026}\label{dat:in:pay-regulation-and-tax-by-location:rates}
\footnotesize
\begin{tabular}{@{}p{2.2cm}p{10.4cm}@{}}
\toprule
London & 20\%, 40\%, 45\% above \pounds125\,140; personal allowance \pounds12\,570, withdrawn above \pounds100\,000; employee
National Insurance 8\% to \pounds50\,270, 2\% above.\\
New York & federal 10\% to 37\% above \$640\,600; state 4\% to 10.9\% (9.65\% from \$1\,077\,550), with the lower rates
recaptured above \$107\,650 of income; city 3.078\% to 3.876\%; Social Security 6.2\% to \$184\,500, Medicare 1.45\% and
0.9\% more above \$200\,000.\\
Chicago & federal and payroll as New York; Illinois 4.95\% flat; no city income tax in the model.\\
Amsterdam & 35.75\% to \euro38\,883, 37.56\% to \euro78\,426, 49.5\% above, national insurance included; credits
withdrawn at higher incomes; expat scheme: up to 30\% of wages tax-free, capped at \euro262\,000 of wages.\\
Zurich & federal tax to 11.5\% of all income at the top; cantonal simple tax to 13\%, multiplied by 2.14 (canton 95\%,
city 119\%); social insurance 5.3\%, and 1.1\% to CHF 148\,200.\\
Singapore & residents 0\% to 24\% above SGD 1\,000\,000.\\
Hong Kong & the lower of 2\% to 17\% on income after allowances and 15\% of income (16\% above HKD 5\,000\,000);
MPF 5\% to HKD 18\,000 a year.\\
Dubai & no income tax on individuals.\\
\bottomrule
\end{tabular}

\smallskip\noindent UK 2026/27; US federal and payroll 2026; New York State and City from the 2025 instructions; the
Netherlands and Zurich 2026; Singapore from YA 2024; Hong Kong 2026/27. Every parameter with its source in
\texttt{data/industry/tax\_2026.csv}.
\end{dated}

\Cref{fig:in:pay-regulation-and-tax-by-location:avg} shows the average rate from \$150\,000 to \$3 million. The places
separate into three groups: Dubai at zero; Hong Kong and Singapore between 10\% and 24\%, Hong Kong nearly flat because
its standard rate caps the average; and the other five between 23\% and 52\%, rising toward their top marginal rates.
New York crosses London near \$1 million, when the state's 9.65\% bracket and its recapture take effect; Zurich starts
low and ends near Chicago.

\begin{omfigure}
\begin{tikzpicture}
\begin{axis}[width=9.2cm,height=6.2cm,xmode=log,xmin=140,xmax=3200,ymin=-2,ymax=56,log ticks with fixed point,
  xtick={150,300,500,1000,2000,3000},xlabel={\small gross pay, \$ thousand (log scale)},ylabel={\small average rate, \%},
  tick label style={font=\footnotesize},grid=major,grid style={gray!20},table/col sep=comma,
  legend style={legend cell align=left,font=\scriptsize,at={(1.02,0.5)},anchor=west}]
\addplot[omDef,thick] table[x=gross,y=new_york]{figdata/industry/15-pay-regulation-and-tax-by-location/avg_rate.csv};
\addplot[omThm,thick] table[x=gross,y=amsterdam]{figdata/industry/15-pay-regulation-and-tax-by-location/avg_rate.csv};
\addplot[black,thick] table[x=gross,y=london]{figdata/industry/15-pay-regulation-and-tax-by-location/avg_rate.csv};
\addplot[omDef,thick,dashed] table[x=gross,y=chicago]{figdata/industry/15-pay-regulation-and-tax-by-location/avg_rate.csv};
\addplot[omThm,thick,dashed] table[x=gross,y=zurich]{figdata/industry/15-pay-regulation-and-tax-by-location/avg_rate.csv};
\addplot[gray,thick] table[x=gross,y=singapore]{figdata/industry/15-pay-regulation-and-tax-by-location/avg_rate.csv};
\addplot[gray,thick,dashed] table[x=gross,y=hong_kong]{figdata/industry/15-pay-regulation-and-tax-by-location/avg_rate.csv};
\addplot[black,thick,dotted] table[x=gross,y=dubai]{figdata/industry/15-pay-regulation-and-tax-by-location/avg_rate.csv};
\addplot[omThm,thick,dotted] table[x=gross,y=amsterdam_expat]{figdata/industry/15-pay-regulation-and-tax-by-location/avg_rate.csv};
\legend{New York,Amsterdam,London,Chicago,Zurich,Singapore,Hong Kong,Dubai,Amsterdam (expat)}
\end{axis}
\end{tikzpicture}

\omcaption{Average rate of income tax and employee contributions on a package of \$150\,000 to \$3 million, a single
employee with no other income, converted at 2025 average exchange rates. Data: \texttt{firm.aftertax} on the 2026
parameters, through \texttt{in\_tax.curve}.}\label{fig:in:pay-regulation-and-tax-by-location:avg}
\end{omfigure}

\section{Expatriate regimes}

\begin{definition}[Tax residence, expatriate tax regime]\label{def:in:pay-regulation-and-tax-by-location:expat}
\emph{Tax residence}\index{tax residence} is the status that makes a person taxable in a country on their income,
usually on worldwide income, decided by that country's rules on days present, home and ties. An \emph{expatriate tax
regime}\index{expatriate tax regime} is a rule that taxes a newly arrived resident more lightly for a limited period,
by exempting a share of pay or foreign income, to attract people with scarce skills.
\end{definition}

Two regimes illustrate the range.
\begin{itemize}
\item \textbf{The Netherlands' expat scheme} lets an employer pay an eligible incomer up to 30\% of wages tax-free for up to
five years, falling to 27\% from 1 January 2027; the tax-free amount is capped by applying the share to at most
\euro262\,000 of wages, a maximum of \euro78\,600. On the one-million package it raises net pay from \$516\,376 to
\$560\,341; its average rate falls from 48.4\% to 44.0\%, while the marginal rate on the bonus stays 49.5\%, because above
the cap every extra euro is taxed in full.
\item \textbf{The United Kingdom} replaced its remittance basis, based on domicile, with ``a new tax regime based on
residence from 6 April 2025'': full relief ``on foreign income and gains for new arrivals \dots{} in their first 4 years of
\omterm{def:in:pay-regulation-and-tax-by-location:expat}{tax residence}'' after ten years away. It relieves income and gains from abroad; it does not by itself reduce the tax on
pay for work done in London.
\end{itemize}
Other countries have regimes of their own, with conditions on qualifications, previous residence and duration; each is
a statute to read, not a rate to assume.

\section{Deferred and equity pay: when it is taxed}

In the United States, property subject to a substantial risk of \omterm{def:in:how-pay-works:rsu}{forfeiture} is income when it vests, at its value then
(chapter~13): deferred pay is taxed when the employee receives it, not when it is awarded. The chapter's model applies the
same timing everywhere; other systems have their own rules for share awards, to be read before a move. Three
consequences follow for a package where the rule applies.
\begin{enumerate}
\item \textbf{The rate is the rate of the vesting year.} An award granted in a year of low income and vesting in a year of
high income is taxed at the high year's marginal rate.
\item \textbf{The value is the value at vesting.} Shares that rise between award and vesting are taxed on the higher value;
shares that fall are taxed on less, but the employee bears the fall (chapter~10 showed the extreme case of tokens).
\item \textbf{A move between countries during vesting} can leave more than one country with a claim on the same award;
the chapter does not model the apportionment, which depends on each country's rules and treaties.
\end{enumerate}

\section{What a number is worth after tax}

The marginal rate is what a bonus meets; \cref{fig:in:pay-regulation-and-tax-by-location:marginal} shows it across the
range of packages. It reaches its ceiling early in London (47\% once the personal allowance has gone) and Amsterdam
(49.5\%), steps up in New York when the state's recapture takes effect (in the model, which approximates the state's worksheets,
the extra tax exceeds the extra pay over a band of \$50\,000 of income above \$1\,077\,550 of taxable income), and stays at or below 24\% in Singapore and 17\% in
Hong Kong throughout. At the low end, withdrawals of allowances and credits push the marginal rate above the top
statutory rate: 62\% in London and 56\% in Amsterdam at \$150\,000.

\begin{center}\footnotesize
\begin{tabular}{@{}lrrrr@{}}
\toprule
\$1 million package (30\% base, 70\% bonus) & local currency & net pay (\$) & average rate & marginal rate\\
\midrule
Dubai & USD 1\,000\,000 & 1\,000\,000 & 0.0\% & 0.0\%\\
Hong Kong & HKD 7\,796\,944 & 844\,104 & 15.6\% & 16.0\%\\
Singapore & SGD 1\,305\,833 & 791\,283 & 20.9\% & 24.0\%\\
Zurich & CHF 829\,243 & 599\,515 & 40.0\% & 44.4\%\\
Chicago & USD 1\,000\,000 & 597\,361 & 40.3\% & 44.3\%\\
London & GBP 758\,235 & 545\,545 & 45.4\% & 47.0\%\\
New York & USD 1\,000\,000 & 540\,584 & 45.9\% & 50.1\%\\
Amsterdam & EUR 884\,969 & 516\,376 & 48.4\% & 49.5\%\\
\bottomrule
\end{tabular}

\smallskip Income tax and employee contributions for a single employee with no other income; 2025 average exchange
rates; the marginal rate is on the last thousand dollars of bonus.
\end{center}

Net pay from the same million ranges from \$516\,376 in Amsterdam to all of it in Dubai, a ratio of 1.94
(\cref{fig:in:pay-regulation-and-tax-by-location:million}). Among the five places whose average rate is 40\% or more,
the spread is narrower: \$516\,376 to \$599\,515, 16\%. The split between base and bonus does not change these figures, because both are
taxed as employment income; it changes the risk (chapter~13) and, under a \omterm{def:in:pay-regulation-and-tax-by-location:cap}{bonus cap}, what the firm may offer.

\begin{omfigure}
\begin{tikzpicture}
\begin{axis}[width=9.2cm,height=6.8cm,xmode=log,xmin=140,xmax=3200,ymin=-3,ymax=120,log ticks with fixed point,
  xtick={150,300,500,1000,2000,3000},xlabel={\small gross pay, \$ thousand (log scale)},
  ylabel={\small marginal rate on the last \$5\,000, \%},tick label style={font=\footnotesize},grid=major,
  grid style={gray!20},table/col sep=comma,legend style={legend cell align=left,font=\scriptsize,at={(1.02,0.5)},anchor=west}]
\addplot[omDef,thick] table[x=gross,y=new_york]{figdata/industry/15-pay-regulation-and-tax-by-location/marginal.csv};
\addplot[omThm,thick] table[x=gross,y=amsterdam]{figdata/industry/15-pay-regulation-and-tax-by-location/marginal.csv};
\addplot[black,thick] table[x=gross,y=london]{figdata/industry/15-pay-regulation-and-tax-by-location/marginal.csv};
\addplot[omDef,thick,dashed] table[x=gross,y=chicago]{figdata/industry/15-pay-regulation-and-tax-by-location/marginal.csv};
\addplot[omThm,thick,dashed] table[x=gross,y=zurich]{figdata/industry/15-pay-regulation-and-tax-by-location/marginal.csv};
\addplot[gray,thick] table[x=gross,y=singapore]{figdata/industry/15-pay-regulation-and-tax-by-location/marginal.csv};
\addplot[gray,thick,dashed] table[x=gross,y=hong_kong]{figdata/industry/15-pay-regulation-and-tax-by-location/marginal.csv};
\addplot[black,thick,dotted] table[x=gross,y=dubai]{figdata/industry/15-pay-regulation-and-tax-by-location/marginal.csv};
\legend{New York,Amsterdam,London,Chicago,Zurich,Singapore,Hong Kong,Dubai}
\end{axis}
\end{tikzpicture}

\omcaption{Marginal rate of income tax and employee contributions on the last \$5\,000 of pay, by gross package, 2026.
The spikes are real features of the schedules: the UK's withdrawal of the personal allowance between \pounds100\,000 and
\pounds125\,140, the Dutch credits' withdrawal and New York's recapture phase-ins. Data: \texttt{in\_tax.marginal\_curve}.}\label{fig:in:pay-regulation-and-tax-by-location:marginal}
\end{omfigure}

\begin{omfigure}
\begin{tikzpicture}
\begin{axis}[width=10.5cm,height=6.4cm,xbar stacked,bar width=9pt,xmin=0,xmax=1080,ymin=0.4,ymax=9.6,ytick=data,
  yticklabels from table={figdata/industry/15-pay-regulation-and-tax-by-location/million.csv}{place},
  yticklabel style={font=\footnotesize},xticklabel style={font=\footnotesize},xlabel={\small \$ thousand},
  grid=major,grid style={gray!20},table/col sep=comma,legend style={legend cell align=left,font=\scriptsize,at={(0.5,-0.2)},anchor=north,
  /tikz/every even column/.append style={column sep=8pt}},legend columns=3]
\addplot[fill=omDef!55,draw=omDef] table[x=net,y=pos]{figdata/industry/15-pay-regulation-and-tax-by-location/million.csv};
\addplot[fill=omThm!45,draw=omThm] table[x=tax,y=pos]{figdata/industry/15-pay-regulation-and-tax-by-location/million.csv};
\addplot[fill=gray!35,draw=gray] table[x=social,y=pos]{figdata/industry/15-pay-regulation-and-tax-by-location/million.csv};
\legend{net pay,income tax,employee contributions}
\end{axis}
\end{tikzpicture}

\omcaption{Where a one-million-dollar package goes, by place of residence and work, in 2026. The Netherlands' expat scheme
is shown as a separate line. Data: \texttt{in\_tax.million}.}\label{fig:in:pay-regulation-and-tax-by-location:million}
\end{omfigure}

The comparison is in dollars, and exchange rates move. At the 2024 average rate a pound bought \$1.2785; at the 2025
average, \$1.3189. A London package fixed in pounds was worth 3.2\% more in dollars in 2025 than in 2024 with no change in
pay, and a comparison made in one year's rates can reverse in the next. A package is paid in the local currency; a
candidate who expects to spend or save in another currency carries that exchange risk.

Tax is not the only difference of place. Chapter~27 sets these figures against the cost of living, the size of each
place's industry and the commute; and whether a given employer has a desk in a given place is a separate question from
what the place's tax would leave.

\section{Tutorial: a million, eight ways}\label{tut:in:pay-regulation-and-tax-by-location:tut}

\begin{tutorial}
\textbf{Goal.} Compute net pay and the average and marginal rates of packages from \$150\,000 to \$3 million in eight
places. \textbf{End state:} the dated table,
\cref{fig:in:pay-regulation-and-tax-by-location:avg,fig:in:pay-regulation-and-tax-by-location:million} and the table above.
\begin{tutsteps}
\item \textbf{Parameters.} \texttt{data/industry/tax\_2026.csv} holds every bracket, rate, cap and credit with a ledger id;
\texttt{ch\_federal\_tax\_2026.csv} the Swiss federal table's break points.
\item \textbf{Structure.} \texttt{firm.aftertax.net\_local(params, place, gross)} applies each system's structure; New York's
benefit recapture is approximated (\cref{lst:in:pay-regulation-and-tax-by-location:ny}).
\omcode{code/firm/aftertax/firm_aftertax.py}{94}{112}{New York State tax with the recapture of the lower brackets phased
in over \$50\,000, continuous at every bracket edge.}\label{lst:in:pay-regulation-and-tax-by-location:ny}
\item \textbf{Convert and compare.} \texttt{net\_usd} converts at the 2025 average rates; \texttt{marginal} differences
net pay over the last thousand dollars.
\end{tutsteps}
\end{tutorial}

\textbf{What to change next.} Add the occupational pension contribution that Zurich deducts from taxable income; model a
UK bonus paid in the following tax year; try a married couple with one income, which changes the answer in several places.

\section{Build: the after-tax calculator}\label{bld:in:pay-regulation-and-tax-by-location:aftertax}

\begin{build}
\bfield{Purpose} Put every pay figure of the role chapters and of chapter~30's choice on an after-tax basis.
\bfield{Interface} \texttt{firm.aftertax}: \texttt{load(path, fed\_path)}; \texttt{progressive(x, brackets)};
\texttt{net\_local(params, place, gross, expat)}; \texttt{local\_per\_usd}; \texttt{net\_usd}; \texttt{marginal};
\texttt{LOCATIONS}. Python standard library.
\bfield{Rules} Parameters are data with a source each, never literals in the code; one employee, single, no other income,
standard deductions only; bonus and salary taxed alike; simplifications named in the module's documentation.
\bfield{Acceptance tests} \texttt{code/firm/aftertax/tests/}: a UK case at \pounds100\,000 and \pounds200\,000 by hand; the
US federal tax on \$1 million of taxable income; New York's tax continuous and flat at 6.85\% and 9.65\% where the recapture
is complete; Hong Kong's progressive and standard tax equal at the official threshold of HKD 2\,132\,500; Singapore's
table values; the Swiss federal table's printed values, including the rounding to CHF 100; the Dutch credits and expat cap.
\bfield{Stretch} Married and family cases; pension deductions; the timing of bonus withholding; a second country during
vesting.
\end{build}

\begin{omsources}
\item HMRC, Rates and thresholds for employers 2026 to 2027; gov.uk, Income Tax rates and Personal Allowances; HM Treasury
and HMRC, Reforming the taxation of non-UK domiciled individuals.
\item IRS, inflation adjustments for 2026 and Topic 751; SSA, contribution and benefit base; New York State, IT-201-I
(2025) and tax expenditure report; Illinois Department of Revenue.
\item Belastingdienst, loonheffingen 2026 tables; Business.gov.nl, the expat scheme.
\item ESTV, Kantonsblatt Zürich and the 2026 federal tax table; Stadt Zürich, Steuerfuss 2026; AHV/IV leaflets 2.01 and
2.08.
\item IRAS, individual income tax rates; GovHK, salaries tax rates; IRD, PAM 61; MPFA; UAE Government portal.
\item Directive 2013/36/EU, Article 94; European Banking Authority, high earners 2024.
\end{omsources}

\section{Exercises}

\begin{exercise}[$\star$]\label{exo:in:pay-regulation-and-tax-by-location:1}
Compute the UK income tax and employee National Insurance on \pounds100\,000 of salary in 2026/27.
\end{exercise}

\begin{exercise}[$\star$]\label{exo:in:pay-regulation-and-tax-by-location:2}
A bank risk-taker in the EU has fixed pay of \euro400\,000. What is the most variable pay the cap allows, with and without
the owners' approval?
\end{exercise}

\begin{exercise}[$\star$]\label{exo:in:pay-regulation-and-tax-by-location:3}
From the table, which place has the highest marginal rate on the bonus of a \$1 million package, and what is it?
\end{exercise}

\begin{exercise}[$\star\star$]\label{exo:in:pay-regulation-and-tax-by-location:4}
Show that Hong Kong's progressive and standard-rate tax are equal at HKD 2\,132\,500 with the basic allowance of HKD
145\,000.
\end{exercise}

\begin{exercise}[$\star\star$]\label{exo:in:pay-regulation-and-tax-by-location:5}
Why does the Dutch expat scheme lower the average rate on a \$1 million package but not its marginal rate?
\end{exercise}

\begin{exercise}[$\star\star$]\label{exo:in:pay-regulation-and-tax-by-location:6}
Why does New York's average rate rise from 45.9\% at \$1 million to 50.1\% at \$1.5 million, when the federal and city rates
are unchanged over that range?
\end{exercise}

\begin{exercise}[$\star\star\star$]\label{exo:in:pay-regulation-and-tax-by-location:7}
\emph{Coding.} With \texttt{firm.aftertax}, find, to the nearest thousand dollars, the gross pay at which New York's net
pay first falls below London's, and explain the crossing.
\end{exercise}

\begin{exercise}[$\star\star\star$]\label{exo:in:pay-regulation-and-tax-by-location:8}
\emph{Find the flaw.} ``Dubai pays twice as much: a million there is worth the same as two million in Amsterdam.''
\end{exercise}

\section{Problem: A Million, Eight Ways}

\begin{problem}[{Weekend problem --- a million, eight ways}]\label{pb:in:pay-regulation-and-tax-by-location:1}
A trader is offered the same package, one million dollars a year with 30\% base and 70\% bonus, in eight offices of one
firm. What does each leave, and what else should decide?

\medskip\noindent\textbf{Part I --- The rules.}
\begin{enumerate}
\item Define a \omterm{def:in:pay-regulation-and-tax-by-location:cap}{bonus cap} and state the EU's.
\item Whom does the cap cover, and whom not, since 2021 and 2023?
\item What does a cap do to the mix and to the level of pay?
\item Define marginal and \omterm{def:in:pay-regulation-and-tax-by-location:rates}{average tax rates}.
\item Does the split between base and bonus change the year's tax in these systems? Why?
\end{enumerate}

\medskip\noindent\textbf{Part II --- The systems.}
\begin{enumerate}[resume]
\item State the UK's rates, allowance and National Insurance for 2026/27.
\item State the US federal, New York State and City rates and payroll taxes.
\item State the Dutch box 1 rates and the expat scheme's terms.
\item State Zurich's structure: federal, cantonal, city, social.
\item State Singapore's, Hong Kong's and Dubai's.
\end{enumerate}

\medskip\noindent\textbf{Part III --- The numbers.}
\begin{enumerate}[resume]
\item Give the package in each local currency at 2025 rates.
\item Give net pay, average and marginal rates in each place.
\item How much does the Dutch expat scheme add?
\item Define \omterm{def:in:pay-regulation-and-tax-by-location:expat}{tax residence} and an expatriate regime; what does the UK's 2025 regime relieve?
\item When is deferred pay taxed, and what three consequences follow?
\end{enumerate}

\medskip\noindent\textbf{Part IV --- The verdict.}
\begin{enumerate}[resume]
\item State the \emph{named result}: net pay from the package in each of the eight places, and the ratio of the highest to
the lowest.
\item Among the five places whose average rate is 40\% or more, what is the spread?
\item What does the model leave out that could change the ranking for her?
\item Why is the tax ranking not the ranking of where to work?
\item In two sentences, advise her.
\end{enumerate}
\end{problem}

\section{Interview questions}

\begin{interviewq}[$\star$ \iqroles{developer}]\label{iq:in:pay-regulation-and-tax-by-location:1}
Write a function for tax under a progressive schedule given as brackets. How do you test it?
\end{interviewq}

\begin{interviewq}[$\star$ \iqroles{bank, risk}]\label{iq:in:pay-regulation-and-tax-by-location:2}
Why did the UK remove the \omterm{def:in:pay-regulation-and-tax-by-location:cap}{bonus cap}, and what did the EU keep?
\end{interviewq}

\begin{interviewq}[$\star\star$ \iqroles{researcher}]\label{iq:in:pay-regulation-and-tax-by-location:3}
An \omterm{def:in:pay-regulation-and-tax-by-location:rates}{average tax rate} curve flattens in one place and rises in another. What features of the two systems explain each?
\end{interviewq}

\begin{interviewq}[$\star\star$ \iqroles{trader}]\label{iq:in:pay-regulation-and-tax-by-location:4}
Your bonus is 70\% of your pay. Would you rather it be taxed at the marginal rate of the year it is awarded or the year it
vests? When does the answer change?
\end{interviewq}

\begin{interviewq}[$\star\star$ \iqroles{developer, mle}]\label{iq:in:pay-regulation-and-tax-by-location:5}
Tax parameters change every year. How would you store them so that last year's calculation can be reproduced exactly?
\end{interviewq}

\begin{interviewq}[$\star\star\star$ \iqroles{researcher, risk}]\label{iq:in:pay-regulation-and-tax-by-location:6}
A model of after-tax pay leaves out pensions, allowances and family status. How would you bound the error for a single
\omterm{def:in:pay-levels-by-role-firm-type-and-seniority:median}{high earner}, and where is it largest?
\end{interviewq}
